% Options for packages loaded elsewhere
\PassOptionsToPackage{unicode}{hyperref}
\PassOptionsToPackage{hyphens}{url}
\documentclass[
  a4paper,
]{article}
\usepackage{xcolor}
\usepackage[margin=25mm]{geometry}
\usepackage{amsmath,amssymb}
\setcounter{secnumdepth}{-\maxdimen} % remove section numbering
\usepackage{iftex}
\ifPDFTeX
  \usepackage[T1]{fontenc}
  \usepackage[utf8]{inputenc}
  \usepackage{textcomp} % provide euro and other symbols
\else % if luatex or xetex
  \usepackage{unicode-math} % this also loads fontspec
  \defaultfontfeatures{Scale=MatchLowercase}
  \defaultfontfeatures[\rmfamily]{Ligatures=TeX,Scale=1}
\fi
\usepackage{lmodern}
\ifPDFTeX\else
  % xetex/luatex font selection
\fi
% Use upquote if available, for straight quotes in verbatim environments
\IfFileExists{upquote.sty}{\usepackage{upquote}}{}
\IfFileExists{microtype.sty}{% use microtype if available
  \usepackage[]{microtype}
  \UseMicrotypeSet[protrusion]{basicmath} % disable protrusion for tt fonts
}{}
\makeatletter
\@ifundefined{KOMAClassName}{% if non-KOMA class
  \IfFileExists{parskip.sty}{%
    \usepackage{parskip}
  }{% else
    \setlength{\parindent}{0pt}
    \setlength{\parskip}{6pt plus 2pt minus 1pt}}
}{% if KOMA class
  \KOMAoptions{parskip=half}}
\makeatother
\usepackage{longtable,booktabs,array}
\newcounter{none} % for unnumbered tables
\usepackage{calc} % for calculating minipage widths
% Correct order of tables after \paragraph or \subparagraph
\usepackage{etoolbox}
\makeatletter
\patchcmd\longtable{\par}{\if@noskipsec\mbox{}\fi\par}{}{}
\makeatother
% Allow footnotes in longtable head/foot
\IfFileExists{footnotehyper.sty}{\usepackage{footnotehyper}}{\usepackage{footnote}}
\makesavenoteenv{longtable}
\usepackage{graphicx}
\makeatletter
\newsavebox\pandoc@box
\newcommand*\pandocbounded[1]{% scales image to fit in text height/width
  \sbox\pandoc@box{#1}%
  \Gscale@div\@tempa{\textheight}{\dimexpr\ht\pandoc@box+\dp\pandoc@box\relax}%
  \Gscale@div\@tempb{\linewidth}{\wd\pandoc@box}%
  \ifdim\@tempb\p@<\@tempa\p@\let\@tempa\@tempb\fi% select the smaller of both
  \ifdim\@tempa\p@<\p@\scalebox{\@tempa}{\usebox\pandoc@box}%
  \else\usebox{\pandoc@box}%
  \fi%
}
% Set default figure placement to htbp
\def\fps@figure{htbp}
\makeatother
\setlength{\emergencystretch}{3em} % prevent overfull lines
\providecommand{\tightlist}{%
  \setlength{\itemsep}{0pt}\setlength{\parskip}{0pt}}
\usepackage{bookmark}
\IfFileExists{xurl.sty}{\usepackage{xurl}}{} % add URL line breaks if available
\urlstyle{same}
\hypersetup{
  hidelinks,
  pdfcreator={LaTeX via pandoc}}

\author{}
\date{}

\begin{document}

\section{Genesis of the Three Spatial Axes and Proper Time:
Re-Derivation of Readouts from the Condensate, and the Uniquely
Unreadable Coordinate
Time}\label{genesis-of-the-three-spatial-axes-and-proper-time-re-derivation-of-readouts-from-the-condensate-and-the-uniquely-unreadable-coordinate-time}

\textbf{Author:} Noriaki Kihara (WF System Co., Ltd.) \textbf{Date:}
2026-08-06 \textbf{Version:} v1 \textbf{Version DOI:}
\href{https://doi.org/10.5281/zenodo.21816652}{10.5281/zenodo.21816652}
\textbf{Concept DOI:}
\href{https://doi.org/10.5281/zenodo.21816651}{10.5281/zenodo.21816651}

\begin{center}\rule{0.5\linewidth}{0.5pt}\end{center}

\subsection{Abstract}\label{abstract}

The preceding paper demonstrated shape-invariant translation of matter
packets, but supplied the position axes (a three-dimensional register
lattice) and the clock (harmonic-closure dispersion) externally. This
paper \textbf{re-derives both} from the observables of the
three-direction metastable condensate itself, re-derives the readouts of
the mass-like and charge-like quantities in the same framework, and
states every readout method in the text. Measured results: (1) the
substance of the condensate is not three equal axes but \textbf{2+1} ---
the only independent substance is one complex eigenpair (a single
rotation plane), and the third direction is dependent, being the
rotation axis of the generator (agreement with the cross product 0.987
\(\to\) 0.9965, improving with N). (2) The spatial \(x, y\) axes are
derived as the \textbf{quadratures} of the plane's complex coordinate
demodulated by the proper clock --- orthogonality is a construction of
the readout, not a property of the substance. (3) Proper time \(\tau\)
is the tick of the condensate's \textbf{single collective clock}
(\(\omega\) \(\approx\) \(\pi\)/72 per step; all occupied planes rotate
rigidly), and this clock is collectively enforced by
\textbf{entrainment}: clockless coherent perturbations are synchronized
to the condensate clock and brought to rest. Displacements of the rest
component persist undamped, and motion becomes the privilege of content
carrying its own clock --- mass. (4) Mass reads as the
\textbf{non-coherence}
\(\det\Gamma = T^2\!-\!X^2\!-\!Y^2\!-\!Z^2 \ge 0\) of a readout pair
(the mathematics is known from the coherence-matrix theory of
polarization; our contribution is the physical identification and the
in-model measurement: light-like eigenmode \(1.2\times10^{-7}\)
vs.~condensate \(3.1\times10^{-3}\)). (5) Charge separates via
transmission, relative phase, and winding number --- readouts of
waveform structure. (6) The \textbf{distributions} of the derived
readouts are measured: sea modes concentrate at maximal non-coherence in
both the growth and the thermal epoch --- the sea is ``heavy''; charge
(winding) distributes over multiple quantized values and does not
stabilize at \(\pm 1\) --- recorded as a known open problem; the
fermionic content of the symmetric sea is gray (particle--antiparticle
balanced), while structures with charged seeds concentrate on the
particle side; matter-packet motion is visualized in the
four-dimensional xyz+\(\tau\) readout as uniform translation (all 10
points at deviation 0.000 for \(\tau\) = 0--900). Organizing all of this
on the closure cone \(x^2+y^2+z^2 = t^2+R^2+Q^2 = R'^2\): \textbf{the
full content of observation is the three spatial axes + proper time
\(\tau\) + mass m + charge q, and the uniquely unreadable quantity is
the coordinate time t} --- t appears only as the arithmetic residue
\(t^2 = R'^2 - m^2 - q^2\), or as a selection that postulates a
wide-area shared time over the spatial oscillations of a large system.
Consequently, the distinction that general relativity postulates ---
proper time is physical, coordinate time is convention --- is
\textbf{derived} from the readout structure. Dimensional sharpness is
finite (ambiguity \(\sim 1/M\)), and at N=4 the third dimension does not
crystallize --- the few-body composite (baryon) domain is cut out of
this paper as a separate subject.

\begin{center}\rule{0.5\linewidth}{0.5pt}\end{center}

\subsection{1. The Character of This Paper --- Separating Premises,
Derivations, and
Measurements}\label{the-character-of-this-paper-separating-premises-derivations-and-measurements}

This paper proposes no new dynamics. Three layers are stated separately.

\begin{itemize}
\tightlist
\item
  \textbf{(a) Premises}: the inherited dynamics (the linear dynamics of
  Papers 7 and 8, read-only import with recorded SHA) and the
  three-direction metastable condensate (discovered in Paper 7; its
  genesis measured end-to-end in the preceding paper).
\item
  \textbf{(b) Derivations}: the construction of the readout methods for
  \(x, y, z\), proper time \(\tau\), mass m, and charge q (Sections
  6--9). \textbf{Citations are provenance labels, not substitutes for
  derivation} --- every readout is re-derived in the text so that
  verification closes within this paper alone.
\item
  \textbf{(c) Measurements}: the preliminary experiment series v1--v5
  and MC (see the experiment inventory; all result JSONs with contrast
  checks).
\end{itemize}

\subsection{2. Claims}\label{claims}

\textbf{First --- the purpose of this paper --- re-derivation of the
readouts.} The readout methods for the three spatial axes \(x,y,z\),
proper time \(\tau\), the mass-like quantity m, and the charge-like
quantity q are re-derived from the observables of the three-direction
metastable condensate alone, with no hand-placed register lattice and no
hand-placed dispersion clock. Each derivation is stated in the text
(Sections 6--9).

\textbf{Second, the structure of the substance (measured, Section 5).}
The substance of the condensate is not three equal axes but
\textbf{2+1}: the top of the occupancy is one complex eigenpair (a
single rotation plane; the conjugate pair of the tangent map with
\(|\lambda|=0.99998\)), and the third direction is not an independent
degree of freedom but the \textbf{rotation axis} of the antisymmetric
flow obtained by projecting the generator onto the occupied subspace ---
its agreement with the cross product of the top two directions is 0.987
(N=5) \(\to\) 0.9965 (N=8), improving with N.

\textbf{Third, derivation of the three spatial axes (Section 6).} \(x\)
and \(y\) are derived as the quadratures of the plane's complex
coordinate \(c(n)\) demodulated by the proper-clock phase
\(\varphi(n)\): \(x=\mathrm{Re}\,(c\,e^{-i\varphi})\),
\(y=\mathrm{Im}\,(c\,e^{-i\varphi})\). \(z\) is the demodulated
projection onto the rotation axis. \textbf{Orthogonality is a
construction of the readout, not a property of the substance.} The 3D
register the preceding paper supplied externally was an external
surrogate of this intrinsic readout.

\textbf{Fourth, derivation of proper time \(\tau\) (Section 7).}
\(\tau\) is the tick of the R' clock --- the condensate's single
collective clock. Measured: all occupied planes rotate rigidly at the
single frequency \(\omega\) \(\approx\) \(\pi\)/72 per step (agreeing
with the fundamental angle of the dynamics to 0.04\%). Moreover this
clock is collectively enforced by \textbf{entrainment}: clockless
coherent perturbations cannot keep their own detuning and synchronize to
the condensate's clock. Rest = clock lock; displacements of the rest
component persist undamped. Motion is the privilege of content with its
own clock --- mass --- whose demonstration was already given by the
matter-packet translation of the preceding paper. The whole is
visualized in the four-dimensional xyz+\(\tau\) readout as uniform
translation (Figure 8: the equally spaced points of the 3D trajectory
and the straight lines of the x--\(\tau\), y--\(\tau\), z--\(\tau\)
projections; all 10 points at deviation 0.000 for \(\tau\) = 0--900).

\textbf{Fifth, re-derivation of the mass-like quantity m (Section 8).} m
is the magnitude of the R component and reads from waveform
\textbf{structure} in three ways: (i) non-coherence --- the determinant
of the Gram matrix of a readout pair, \(\det\Gamma = T^2-X^2-Y^2-Z^2\)
(fully derived in the text; the mathematics is known from polarization
theory, the identification is this paper's contribution); (ii)
closure-completion cost; (iii) the strength of one's own clock.

\textbf{Sixth, re-derivation of the charge-like quantity q (Section 9).}
q is the magnitude plus relative sign of the Q component, and reads from
waveform structure in three ways: (i) transmission, (ii) relative phase
(the charge grammar), (iii) winding number (Kaluza--Klein-type
momentum).

\textbf{Seventh, distributions of the readouts (measured supplement,
Section 9.5).} Applying the derived readouts to the mode set \((k,m)\)
in three states (the symmetric sea in its growth and thermal epochs; a
charged structure): (a) sea modes concentrate at \textbf{maximal
non-coherence (mass\(^2\)/T\(^2\) \(\approx\) 1)} in both epochs --- the
sea is heavy, and light (light-like) states exist only as coherent
structures. (b) \textbf{Charge (winding) distributes over multiple
quantized values and does not stabilize at \(\pm 1\)} --- the origin of
the uniqueness of the elementary charge is a known open problem, and
this paper provides the direct measurement of the distribution (Section
15-6). (c) In the spin-like quantity (the Bloch vector of the Gram ---
an exploratory readout) the sea is essentially unpolarized. (d)
Fermionic waves classify by sign balance into \textbf{particle-like,
antiparticle-like, and gray}; the symmetric sea concentrates on gray,
structures with a charged seed on the particle side --- the
particle/antiparticle distinction does not exist in the symmetric sea;
it is brought in by charged seeds.

\textbf{Eighth --- the central thesis --- the full content of
observation, and the single unreadable (Section 10).} On the closure
cone \(x^2+y^2+z^2 = t^2+R^2+Q^2 = R'^2\), the full content of
observation is \textbf{the three spatial axes + proper time \(\tau\) +
mass m + charge q}. \textbf{The uniquely unreadable quantity is the
coordinate time t}: no direct readout instrument for t exists; t appears
only as the arithmetic residue \(t^2=R'^2-m^2-q^2\), or as the
\textbf{selection} that postulates a wide-area shared time over the
spatial oscillations of a large system. Two consequences. (1) The
distinction general relativity postulates --- proper time is physical,
coordinate time is convention --- is \textbf{derived} from the readout
structure. (2) A candidate answer to why observed spacetime is 3+1 ---
because the readable content is exactly three real axes plus one clock.

\textbf{Ninth, conventions of orientation (Section 11).} There is no
arrow of time. The imaginary sign is the bookkeeping of ``the direction
of this axis cannot be read'' --- the closure sees only squares.
Physical quantities are relative signs only; the relative handedness of
clock \(\times\) rotation axis (-) is measured stably as a
convention-invariant.

\textbf{Tenth, the boundary of applicability (Section 12).} The
sharpness of the third dimension is finite: ambiguity
\(1-|\hat n\cdot\hat a| \sim 1/M\). At N=4 the third dimension does not
crystallize. N=4,5 is the domain of few-body composites (baryons), where
dimension, time, and even existence are ambiguous --- that is the
physics of that domain; this paper's subject is the large-N side where
space exists, and the few-body domain is cut out as a separate subject.

\textbf{Eleventh, the full record of refutations and corrections
(Section 13)} --- five instrument errors, one of which is a live example
of ``the limits of resolution fabricating structure'': the paper's own
subject appearing on the instrument side.

\textbf{Finally, what this paper does not claim, and the falsification
conditions (Section 16).}

\subsection{3. Motivation --- the Two Things Supplied
Externally}\label{motivation-the-two-things-supplied-externally}

The preceding paper (the many-body connection paper {[}1{]}) ran vacuum
\(\to\) inflation \(\to\) condensation \(\to\) matter precipitation
under a single dynamics with no timing rule, and demonstrated that the
generated matter packets move at constant velocity without changing
shape. But that motion demo had two hand-placed inputs: \textbf{the
position axes} (a \(8^3\) register lattice) and \textbf{the clock} (the
harmonic-closure dispersion \(\omega_k = k\,\omega_1\) {[}6{]} inserted
as a register clock).

With space and time supplied externally, one has not yet earned the
title ``genesis of the three spatial axes and proper time.'' The
question of this paper is: \textbf{can the condensate supply the
position axes and the clock from its own observables?} The answer is
affirmative --- and stronger than expected: the three axes are not equal
but 2+1, orthogonality is manufactured by the readout, the clock is
single and collectively enforced by entrainment, and exactly one
quantity remains unreadable.

\subsection{4. Premises --- Dynamics and the
Condensate}\label{premises-dynamics-and-the-condensate}

The dynamics is inherited unmodified from Papers 7 and 8 {[}3,4{]}
(read-only import, recorded SHA series). Relational waves
\(Z \in \mathbb{C}^M\) of \(N\) bodies (\(M=N(N-1)/2\)), generator
\(K_{ef}=\sin(\theta_f-\theta_e)\) (adjacent edges, antisymmetric),
Cayley step. Starting from a self-consistent closure the system reaches,
after latency and burst, the \textbf{three-direction metastable
structure} (the discovery of Paper 7). All measurements of this paper
are performed on this metastable condensate (post-settling windows; see
the experiment inventory). Hereafter the \textbf{update count (the
computational step label) is written \(n\)} --- this is distinct from
the coordinate time \(t\) that appears in the closure cone (\(t\) is the
subject of Section 10 and appears in this paper only inside the
decomposition of the cone). The argument of the parent plane's complex
coordinate (the plane spanned by the real and imaginary parts of the
self-consistent solution \(v\)) is written \(\varphi(n)\), and the
component with the parent plane projected out is written \(Z_\perp\).

\subsection{5. The Structure of the Substance --- 2+1
(Measured)}\label{the-structure-of-the-substance-21-measured}

Three measurements on the metastable-window trajectory pin down the
substance of the condensate (Figure 1).

First, the singular values of the \(Z_\perp\) trajectory come
\textbf{strictly in pairs} (1.000, 0.998 / 0.427, 0.426 / 0.207, 0.207 /
\ldots) --- the occupancy is a hierarchy of rotation planes. Second, in
the eigendecomposition of the one-step tangent matrix (\(2M\times2M\),
finite differences), the occupied subspace is carried by \textbf{one
complex-conjugate eigenpair} (\(|\lambda|=0.999984\)) --- the
independent substance is a single rotation plane, and no independent
eigenvalue exists for a third direction. Third, the rotation axis
\(\hat a\) of the antisymmetric \(3\times3\) obtained by projecting the
generator \(K\) onto the occupied 3D subspace agrees with the cross
product of the plane's two directions,
\(\hat n = \hat d_1\times\hat d_2\): \(|\hat n\cdot\hat a| = 0.9874\)
(N=5) \(\to\) 0.9918 (N=6) \(\to\) 0.9965 (N=8).

\begin{figure}
\centering
\pandocbounded{\includegraphics[keepaspectratio,alt={Figure 1 the 2+1 structure}]{figs/fig_s1_structure_v1.png}}
\caption{Figure 1 the 2+1 structure}
\end{figure}

\textbf{Figure 1}: (a) singular values come strictly in pairs = a
hierarchy of rotation planes. (b) Tangent-map eigenvalues: the occupancy
is carried by one conjugate pair (red). The third direction has no
independent eigenvalue.

In sum, the substance is \textbf{one plane (one complex degree of
freedom) + one dependent rotation axis} --- the substance of the ``three
directions'' is not 3 but 2+1; more precisely, complex-1 + real-1.

\subsection{6. Derivation of the Three Spatial
Axes}\label{derivation-of-the-three-spatial-axes}

\textbf{Derivation.} A state rotating within the plane is perfectly
correlated in its projections onto \textbf{any} two in-plane directions
(measured: readout correlation 1.0000) --- hence the plane does not
possess two independent real readouts. The plane's one independent
readout is the complex coordinate

\[c(n) = \langle d_{\text{plane}},\, Z_\perp(n)\rangle\]

(where \(d_{\text{plane}}\) is the complex direction of the plane
constructed from the eigenpair). This is demodulated by the proper-clock
phase \(\varphi(n)\) (Section 7) and \textbf{decomposed into
quadratures}:

\[x(n) = \mathrm{Re}\,\big(c\,e^{-i\varphi}\big), \qquad y(n) = \mathrm{Im}\,\big(c\,e^{-i\varphi}\big).\]

An important point must be made explicit: \(\varphi(n)\) is \textbf{not}
the argument of \(c(n)\) itself (if it were, \(c\,e^{-i\varphi}=|c|\)
and \(y\equiv0\) --- the readout would degenerate). \(\varphi\) is the
collective clock phase read independently from the parent plane, and
\(c\) is a coordinate of \(Z_\perp\) --- two different readouts --- so
\(c\,e^{-i\varphi}\) gives the \textbf{relative complex coordinate with
respect to the clock}. The orthogonality of \(x\) and \(y\) is the
orthogonality of the real and imaginary parts of a complex number ---
that is, a \textbf{construction of the readout} --- and does not
correspond to two orthogonal axes on the substance side (the substance
is one complex degree of freedom). The third axis is the demodulated
projection onto the rotation axis:

\[z(n) = \mathrm{Re}\,\big(\langle \hat a, Z_\perp\rangle\, e^{-i\varphi}\big).\]

The seemingly isotropic orthogonal three axes are thus manufactured as
``the phase decomposition of one complex coordinate + the dependent
axis.''

\textbf{Measurements} (Figure 2): the clock-synchronized component keeps
a constant radius \(|r|=0.983\pm0.005\) and only precesses slowly in
angle (0.12\% per clock) --- quasi-rest with invariant amplitude. A
displacement (kick) of the rest component persists undamped for 8000
steps.

\begin{figure}
\centering
\pandocbounded{\includegraphics[keepaspectratio,alt={Figure 2 quadrature readout}]{figs/fig_s2_quadrature_v1.png}}
\caption{Figure 2 quadrature readout}
\end{figure}

\textbf{Figure 2}: (a) quasi-rest of the clock-synchronized component.
(b) persistence of a displacement.

The 3D register lattice supplied externally by the preceding paper turns
out to have been an external surrogate of this intrinsic readout.

\subsection{\texorpdfstring{7. Derivation of Proper Time
\(\tau\)}{7. Derivation of Proper Time \textbackslash tau}}\label{derivation-of-proper-time-tau}

\textbf{Derivation.} As the clock corresponding to the right side of the
cone, \(R'^2 = t^2+R^2+Q^2\), what a local system can read is the tick
of the total \(R'\) (the decomposition is Section 10). In the condensate
this appears as the \textbf{rotation phase \(\varphi(n)\) of the parent
plane} --- the single collective clock. The proper-time readout is

\[\tau(n) = \varphi(n)\,/\,\omega_{\text{clock}}.\]

\textbf{Measurement 1: the clock is single (Figure 3a).} At high
resolution (10\(\times\) window, zero padding), the rotation frequencies
of all occupied planes converge to the single value \(\omega\)
\(\approx\) \(\pi\)/72 per step (agreeing with the fundamental angle of
the dynamics to 0.04\%; ratios \(1\pm1\%\)). There is no ladder of
frequencies (the ``ladder'' seen at low resolution was an instrument
artifact --- Section 13).

\textbf{Measurement 2: the clock is collectively enforced by entrainment
(Figure 3b).} The natural fluctuation content of the condensate shows
asynchronous phase with slips, but a \textbf{coherent} kick in the same
direction cannot keep its own detuning and locks to the condensate clock
(drift \(-1.4\% \to -0.05\%\)). Proper time is not an attribute of
individual components; it is \textbf{an order that the condensate
enforces collectively through entrainment}.

\begin{figure}
\centering
\pandocbounded{\includegraphics[keepaspectratio,alt={Figure 3 single clock and entrainment}]{figs/fig_s3_clock_entrainment_v1.png}}
\caption{Figure 3 single clock and entrainment}
\end{figure}

\textbf{Figure 3}: (a) all occupied planes rotate rigidly at a single
clock (four values of N). (b) Entrainment --- perturbations are
synchronized to the condensate's proper clock.

\textbf{Consequence: the classification of rest and motion.} Rest =
clock lock (displacements persist, Figure 2b). All clockless
perturbations are entrained to rest. Hence \textbf{motion is the
privilege of content that keeps ticking its own phase against
entrainment --- that which carries its own clock = mass.} The mass-side
demonstration was already given by the matter-packet translation of the
preceding paper {[}1{]} --- Figure 8 visualizes the whole in this
paper's readout vocabulary: the 3D trajectory of the xyz readout draws
an equally spaced sequence of points (uniform translation), and in all
three projection planes x--\(\tau\), y--\(\tau\), z--\(\tau\), the
measured points lie on the predicted straight lines (\(\tau\) = 0--900,
all 10 points at deviation 0.000, PR constant 3.39 --- the packet
travels straight for 900 steps without changing shape at all). The chain
mass \(\to\) proper clock \(\to\) time \(\to\) motion closes.

\begin{figure}
\centering
\pandocbounded{\includegraphics[keepaspectratio,alt={Figure 8 4D translation}]{figs/fig_s8_translation4d_v1.png}}
\caption{Figure 8 4D translation}
\end{figure}

\textbf{Figure 8}: the four-dimensional xyz+\(\tau\) (R' proper time)
readout. (a) 3D trajectory --- equally spaced points = uniform
translation. (b)(c)(d) projections onto the x--\(\tau\), y--\(\tau\),
z--\(\tau\) planes --- straight lines = constant velocity. All 10 points
at deviation 0.000.

\subsection{8. Re-Derivation of the Mass-Like Quantity
m}\label{re-derivation-of-the-mass-like-quantity-m}

Mass is the magnitude of the R component and reads from waveform
\textbf{structure} (quantities visible in a snapshot). Three routes are
derived.

\subsubsection{8.1 Readout (i): non-coherence (fully
derived)}\label{readout-i-non-coherence-fully-derived}

For a two-component complex readout \(u=(c_1, c_2)\) (e.g., the
parent-plane coordinate and the principal out-of-plane coordinate), the
window-averaged Gram matrix

\[\Gamma = \langle u\,u^\dagger\rangle = \begin{pmatrix} \langle|c_1|^2\rangle & \langle c_1\bar c_2\rangle \\ \langle c_2\bar c_1\rangle & \langle|c_2|^2\rangle \end{pmatrix}\]

is \(2\times2\) Hermitian. Decomposing in the Pauli basis,
\(\Gamma = T\,\mathbb{1} + X\sigma_x + Y\sigma_y + Z\sigma_z\), the
determinant is identically

\[\det\Gamma = T^2 - X^2 - Y^2 - Z^2.\]

Since \(\Gamma\) is positive semidefinite, \(\det\Gamma \ge 0\) ---
\textbf{a light-cone-type bound holds automatically}. Equality
\(\det\Gamma=0\) holds if and only if \(\Gamma\) has rank 1, i.e., \(u\)
is proportional to a single complex amplitude within the window
(\(c_2 = \text{const}\cdot c_1\)) --- fully coherent =
\textbf{light-like}. We therefore take

\[m^2 \equiv \det\Gamma\]

as the mass readout: mass is the \textbf{non-coherence} of the readout
pair. We distinguish normalizations: \(\det\Gamma\) is quartic in
amplitude, so the absolute mass candidate is \(\sqrt{\det\Gamma}\),
while the dimensionless \textbf{mass degree} is
\(\det\Gamma/T^2 \in [0,1]\). What this paper measures (including
Section 9.5) is the latter; ``maximal mass'' means, precisely, maximal
non-coherence degree --- fixing the absolute energy scale (the scale of
the R axis) remains future work.

\textbf{Provenance of the mathematics.} This mathematics is not new: the
Pauli coefficients of a \(2\times2\) coherence matrix are the Stokes
parameters of polarization, and \(\det J \ge 0\) (equality = full
polarization) is a classical identity of polarization theory
{[}11,12{]}. This paper's contribution is the physical
\textbf{identification} --- reading the degree of depolarization as
mass\(^2\) and identifying light-like = fully polarized (rank 1) ---
plus the following in-model measurements.

\textbf{Machine verification} (Figure 4; criteria fixed before
execution): a synthetic rank-1 state gives
\(m^2/T^2 = 1.4\times10^{-16}\) (machine zero); an incoherent synthesis
of two independent components gives \(1.000\). In the model, the
pre-condensation parent eigenmode (the light-like candidate) gives
\(1.2\times10^{-7}\), the metastable condensate \(3.1\times10^{-3}\) ---
\textbf{the condensate carries \(2.5\times10^4\) times the mass degree
of the light-like eigenmode}, consistent with the entrainment of Section
7 (the condensate enforces its own clock).

\begin{figure}
\centering
\pandocbounded{\includegraphics[keepaspectratio,alt={Figure 4 mass = non-coherence}]{figs/fig_s6_mass_check_v1.png}}
\caption{Figure 4 mass = non-coherence}
\end{figure}

\textbf{Figure 4}: the mass readout = non-coherence of the Gram matrix.
Machine verification on four states.

\subsubsection{8.2 Readout (ii): closure-completion
cost}\label{readout-ii-closure-completion-cost}

In the frame reading generation as ``completion of closure'' (Paper 9
{[}7{]}), mass reads as the cost of completing the closure (the
operational definition and the \(f(N)=2\ln N/N\)-type counting are given
in {[}7{]}; the identification is established by that paper's
experiments). Here we keep this as a provenance label; the primary
readouts of this paper are (i) and (iii).

\subsubsection{8.3 Readout (iii): the strength of one's own
clock}\label{readout-iii-the-strength-of-ones-own-clock}

From Section 7, mass reads dynamically as the \textbf{degree of
resistance to entrainment} --- the strength with which content keeps
ticking its own phase. Light-like linear perturbations are entrained
completely (zero mass); matter content keeps its clock and moves.
Consistency with readout (i): fully coherent (rank 1)
\(\Leftrightarrow\) entrained \(\Leftrightarrow\) massless holds
simultaneously in the measurements of Section 7 and Figure 4.

\subsection{9. Re-Derivation of the Charge-Like Quantity
q}\label{re-derivation-of-the-charge-like-quantity-q}

Charge is the magnitude plus relative sign of the Q component and
likewise reads from waveform structure. Three operational definitions
are stated.

\textbf{Readout (i): transmission.} From the two-wave collision readout
\(\theta = \operatorname{atan2}(\sqrt{P_f}, \sqrt{P_b})\) (the ratio of
fermionic to bosonic relational content of the combined spectrum), the
reflection rate \(R=\sin^2\theta\) is read, and its complement --- the
transmission \(1-R\) --- is the elementary-charge readout (the
quantitative identification \(1-R=\sqrt{4\pi\alpha}\) was established in
Trilogy A {[}6{]}).

\textbf{Readout (ii): relative phase (the charge grammar).} The
\textbf{sign} of charge reads as the relative phase of two waves (same
sign = one phase relation, opposite sign = inverted; the grammar
including its two neutralization mechanisms was established in {[}6{]}).
This is consistent with the principle of Section 11 (physical quantities
are relative signs only) --- the sign of charge, too, is relational, not
a standalone attribute.

\textbf{Readout (iii): winding number.} The dominant winding number
\(m\) of the Fourier transform along the internal (\(\eta\)) axis is the
charge-type (Kaluza--Klein) momentum. That pair creation obeys the
winding sum rule \(m^* = 2m_B - m_s\) exactly is measured in the
two-body and many-body censuses {[}2,1{]} (unpredicted windings at
machine zero).

All three read structural quantities from waveform snapshots;
consistently with the closure (which sees only squares), what is
readable is magnitude and relative sign --- not the signed coordinate
itself.

\subsection{9.5 Distributions of the Readouts --- Mass, Charge, Spin,
and Classes in the Sea (Measured
Supplement)}\label{distributions-of-the-readouts-mass-charge-spin-and-classes-in-the-sea-measured-supplement}

The readouts of Sections 8--9 are applied to the mode set \((k, m)\)
(raw bin \(\times\) winding) in three states --- the symmetric sea in
its growth epoch (j=400--500) and its thermal epoch (j=3000--3200), and
a charged structure (census-type: single-winding pumps + a
single-winding seed).

\textbf{The mass distribution (Figure 9)}: sea modes concentrate at
\textbf{maximal non-coherence (mass\(^2\)/T\(^2\) \(\approx\) 1)} in
both epochs --- the sea is ``heavy.'' Light (light-like) states do not
appear in the sea; they are confined to coherent structures (the
eigenmode at \(1.2\times10^{-7}\), Figure 4). The structure of the mass
spectrum is carried by coherent structures, not by the sea.

\textbf{The charge distribution (Figure 10)}: in the symmetric sea,
charge (winding) distributes almost flatly over many values and
\textbf{does not stabilize at \(\pm 1\) --- this is a known open
problem, and this measurement is its direct observation} (Section 15-6).
In the charged structure, multiple charge values coexist: the seed
\(m=1\) dominates and the sum-rule partner \(m=3\) stands clearly above
the machine-zero floor (log display) --- charge takes a \textbf{set} of
quantized values, but the mechanism that collapses it to \(\pm 1\) has
not been found.

\textbf{The distributions of bosonic and fermionic waves, and of the
spin-like quantity (Figure 11)}: the class-resolved power spectra show
the growth-epoch structure (pump lines, the fermionic hump) flattening
under thermalization (class totals F/B 21.9/42.2 \(\to\) 28.9/35.1 ---
approach to equipartition). The spin-like quantity (the Bloch vector of
the Gram --- an \textbf{exploratory} readout, a by-product of Section
8.1) shows the sea essentially unpolarized (\(|S|\approx0.03\)), the
coherent charged seed about five times more polarized
(\(|S|\approx0.13\)) yet far from full polarization. A full
identification of spin is future work.

\textbf{The distribution of particle-like, antiparticle-like, and gray
states (Figure 12)}: the sign-balance \(b=(P_+-P_-)/(P_++P_-)\)
distribution of fermionic waves concentrates near \(b=0\)
(\textbf{gray}) in the symmetric sea (particle and antiparticle
fractions zero in both epochs) and at \(b=+1\) (\textbf{particle-like})
in the structure with a charged seed. Particle/antiparticle/gray is
readable as a three-way classification of states: symmetric initial
conditions grow a gray sea; charged seeds grow particle-like content.

\begin{figure}
\centering
\pandocbounded{\includegraphics[keepaspectratio,alt={Figure 9 mass distribution}]{figs/fig_s9_mass_distribution_v1.png}}
\caption{Figure 9 mass distribution}
\end{figure}

\textbf{Figure 9}: the distribution of the mass readout. The sea
concentrates at maximal non-coherence.

\begin{figure}
\centering
\pandocbounded{\includegraphics[keepaspectratio,alt={Figure 10 charge distribution}]{figs/fig_s10_charge_distribution_v1.png}}
\caption{Figure 10 charge distribution}
\end{figure}

\textbf{Figure 10}: the distribution of the charge readout.
Stabilization at \(\pm 1\) is unresolved (open problem).

\begin{figure}
\centering
\pandocbounded{\includegraphics[keepaspectratio,alt={Figure 11 classes and spin}]{figs/fig_s11_class_spin_v1.png}}
\caption{Figure 11 classes and spin}
\end{figure}

\textbf{Figure 11}: the distributions of bosonic and fermionic waves and
of the spin-like quantity.

\begin{figure}
\centering
\pandocbounded{\includegraphics[keepaspectratio,alt={Figure 12 particle, antiparticle, gray}]{figs/fig_s12_particle_gray_v1.png}}
\caption{Figure 12 particle, antiparticle, gray}
\end{figure}

\textbf{Figure 12}: the distribution of particle-like,
antiparticle-like, and gray states of fermionic waves.

\subsection{10. The Central Thesis --- the Full Content of Observation,
and the Single
Unreadable}\label{the-central-thesis-the-full-content-of-observation-and-the-single-unreadable}

The readouts of Sections 6--9 are organized on the closure cone

\[x^2 + y^2 + z^2 \;=\; t^2 + R^2 + Q^2 \;=\; R'^2\]

(Figure 5).

{\def\LTcaptype{none} % do not increment counter
\begin{longtable}[]{@{}
  >{\raggedright\arraybackslash}p{(\linewidth - 4\tabcolsep) * \real{0.3333}}
  >{\raggedright\arraybackslash}p{(\linewidth - 4\tabcolsep) * \real{0.3333}}
  >{\raggedright\arraybackslash}p{(\linewidth - 4\tabcolsep) * \real{0.3333}}@{}}
\toprule\noalign{}
\begin{minipage}[b]{\linewidth}\raggedright
Quantity
\end{minipage} & \begin{minipage}[b]{\linewidth}\raggedright
Readable?
\end{minipage} & \begin{minipage}[b]{\linewidth}\raggedright
Means (derived in the text)
\end{minipage} \\
\midrule\noalign{}
\endhead
\bottomrule\noalign{}
\endlastfoot
\(x, y, z\) & Yes: direct & quadratures (Section 6) \\
proper time \(\tau\) = the R' clock & Yes: direct & the tick of the
single clock, collectively enforced by entrainment (Section 7) \\
mass m (R component) & Yes: separable & non-coherence / completion cost
/ own-clock strength (Section 8) \\
charge q (Q component) & Yes: separable & transmission / relative phase
/ winding (Section 9) \\
\textbf{coordinate time t} & \textbf{\(\times\) uniquely not} & only as
the residue \(t^2=R'^2-m^2-q^2\) or as a wide-area selection \\
\end{longtable}
}

\begin{figure}
\centering
\pandocbounded{\includegraphics[keepaspectratio,alt={Figure 5 the readout structure of the closure cone}]{figs/fig_s7_cone_readout_v1.png}}
\caption{Figure 5 the readout structure of the closure cone}
\end{figure}

\textbf{Figure 5}: what can be read, and what cannot.

\textbf{The uniquely unreadable quantity is the coordinate time t.}
Surveying the entire experimental corpus of this series, no direct
readout instrument for t exists (the ``time'' of every experiment is
either a step counter = convention, or the R' clock). t appears through
exactly two routes: (a) the arithmetic residue
\(t^2 = R'^2 - m^2 - q^2\); (b) the route that collates the spatial
oscillations of a large system (many local clocks) and \textbf{selects}
the postulate of a wide-area shared time --- indeed, the everyday
separation of t proceeds by this route, and the separations of R and Q
likewise proceed by bookkeeping from spring-like and electromagnetic
oscillations observed \textbf{in space}.

\textbf{A structural reason (candidate)}: the contents of R and Q are
waveform \textbf{structure} (non-coherence, winding --- they appear in a
snapshot), but the content of t is the \textbf{flow} itself, which
leaves no trace in a snapshot. The flow cannot be read from within the
flow.

\textbf{Consequence 1: derivation of the proper-time/coordinate-time
distinction.} In general relativity, ``proper time is what clocks
measure --- physical; coordinate time is gauge --- convention'' follows
from the \textbf{principle} of diffeomorphism invariance. Here the same
distinction is \textbf{derived} from the readout structure of the
closure cone --- the direction of the principle is reversed.

\textbf{Consequence 2: why 3+1.} The structure of observation is exactly
3+1 (three real axes + one clock) plus matter attributes (m, q). Because
the readable content is exactly that --- this is the model's candidate
answer.

\subsection{11. Conventions of Orientation --- There Is No Arrow of
Time}\label{conventions-of-orientation-there-is-no-arrow-of-time}

The model has no arrow of time. The imaginary sign is the bookkeeping of
``the direction of this axis cannot be read'' (\((it)^2=(-it)^2\) ---
the closure sees only squares), and running the map forward is
convention, not physics. Measurements agree: replacing the initial state
by its complex conjugate (conjugation of the state is not time reversal
--- Section 13) leaves the convention-invariant \textbf{relative sign}
--- the relative handedness of the clock rotation and the rotation axis
--- unchanged (-). Physical quantities are relative signs only; absolute
orientations (the arrow of time, the handedness of space) are selection
problems. Note that the claim of this section is that \textbf{the
fundamental closure and the reversible map have no absolute time
orientation}; whether a \textbf{statistical effective arrow of time}
emerges from the history of thermalization, generation, and
autocatalysis is a separate question (Section 15-9).

\subsection{12. The Boundary of Applicability --- Dimensional Sharpness
and the Few-Body
Domain}\label{the-boundary-of-applicability-dimensional-sharpness-and-the-few-body-domain}

The agreement of the third direction's dependence (rotation axis = cross
product) is not perfect and improves with N. The ambiguity
\(1-|\hat n\cdot\hat a|\) scales roughly as the inverse of the number of
relational waves \(M\) (0.0126 \(\to\) 0.0082 \(\to\) 0.0035 for
N=5,6,8), and \textbf{at N=4 the ambiguity is 0.9986 (agreement
\(|\hat n\cdot\hat a|=0.0014\)) --- the third dimension does not
crystallize at all} (Figure 6). Explicitly, the hierarchy is continuous:
N=4 = non-crystallized; N=5 = crystallized but with maximal ambiguity
(the edge of the few-body domain); large N = a sharp spatial phase.

\begin{figure}
\centering
\pandocbounded{\includegraphics[keepaspectratio,alt={Figure 6 dimensional sharpness}]{figs/fig_s4_sharpness_v1.png}}
\caption{Figure 6 dimensional sharpness}
\end{figure}

\textbf{Figure 6}: dimensional ambiguity \textasciitilde1/M. At N=4 the
third dimension does not crystallize.

This is not a defect but the boundary of applicability. As closed
relational systems, N=4,5 is the domain of \textbf{few-body composites
--- baryons} --- extremely quantized and discretized, alias-ridden, with
dimension, time, and even existence ambiguous: that is the physics of
that domain itself (consistent with the separate line of this series
identifying fractional quark charges with low-order rational addresses).
The subject of this paper is the large-N side where space exists; the
few-body domain is explicitly cut out as a separate research subject.

\subsection{13. Record of Refutations and
Corrections}\label{record-of-refutations-and-corrections}

{\def\LTcaptype{none} % do not increment counter
\begin{longtable}[]{@{}
  >{\raggedright\arraybackslash}p{(\linewidth - 6\tabcolsep) * \real{0.1429}}
  >{\raggedright\arraybackslash}p{(\linewidth - 6\tabcolsep) * \real{0.2857}}
  >{\raggedright\arraybackslash}p{(\linewidth - 6\tabcolsep) * \real{0.2857}}
  >{\raggedright\arraybackslash}p{(\linewidth - 6\tabcolsep) * \real{0.2857}}@{}}
\toprule\noalign{}
\begin{minipage}[b]{\linewidth}\raggedright
\#
\end{minipage} & \begin{minipage}[b]{\linewidth}\raggedright
Error
\end{minipage} & \begin{minipage}[b]{\linewidth}\raggedright
Discovery
\end{minipage} & \begin{minipage}[b]{\linewidth}\raggedright
Correction
\end{minipage} \\
\midrule\noalign{}
\endhead
\bottomrule\noalign{}
\endlastfoot
1 & Phase-gradient gauge: frequency measured by the slope of the
demodulated phase & Phase slips at amplitude zero crossings contaminate
it & Switched to FFT-peak gauges. The baseline ``detuning'' numbers
contain slip contamination, so the entrainment conclusion is stated as
``slipping asynchrony \(\to\) slip-free lock'' (Figure 3b) \\
2 & \textbf{FFT bin quantization misread as an ``equally spaced comb at
0.1\% precision''} & The ``spacing'' agreed exactly with the FFT bin
width \(2\pi/n_s\) at every N & Re-measured with a 10\(\times\) window
plus zero padding \(\to\) the comb vanishes; a single clock (Figure 7).
\textbf{The limits of resolution fabricate structure --- a live instance
of this paper's subject appearing on the instrument side} \\
3 & Window transients mistaken for real structure & Vanish at high
resolution in late windows & Standardized on post-settling windows \\
4 & Conjugation of the state confused with time reversal (chirality-test
design error) & The conjugate state also converged to the same sign &
Reinterpreted via ``physical quantities are relative signs only''
(Section 11); the conjugation test stands as a check of convention
invariance \\
5 & The naive design ``motion = detuning'' & Kicks do not keep their
detuning; they lock to the clock & \textbf{The discovery of entrainment}
(Section 7) --- a refutation turned into the central result \\
\end{longtable}
}

\begin{figure}
\centering
\pandocbounded{\includegraphics[keepaspectratio,alt={Figure 7 the false comb}]{figs/fig_s5_false_comb_v1.png}}
\caption{Figure 7 the false comb}
\end{figure}

\textbf{Figure 7}: instrument lesson \#2. The low-resolution ``equally
spaced comb'' (left) is exactly the FFT bin width and vanishes at high
resolution (right).

Three of the five (\#1--2, \#5 with \#4) are the same mistake: building
an instrument without deriving the properties of the measured quantity.
The discipline is written out: \textbf{derive the invariance (or the
predicted value) of the measured quantity before designing the
instrument.}

\subsection{14. Prior Work --- Correspondences Found After Independent
Derivation}\label{prior-work-correspondences-found-after-independent-derivation}

The readout constructions of this paper were derived independently from
the closure cone and the condensate's observables. Re-surveying
afterwards, we find the following correspondences --- supporting
evidence, not proofs.

\begin{enumerate}
\def\labelenumi{\arabic{enumi}.}
\tightlist
\item
  \textbf{Relational time --- the Page--Wootters mechanism} {[}8{]}.
  Time emerges for subsystems from correlations with a clock subsystem
  inside a globally stationary state (``evolution without evolution'').
  Differences: there, conditioning on quantum constraints and
  entanglement; here, the clock itself is dynamically self-generated
  (condensation + entrainment), and the unreadability of coordinate time
  is derived from the readout structure.
\item
  \textbf{Coordinate time is gauge --- partial observables} {[}9{]}.
  That coordinate time is not an observable in general relativity is
  orthodox, formalized by Rovelli's partial observables. Difference: in
  GR the distinction follows from the principle of diffeomorphism
  invariance; here it is derived from the readout structure ---
  \textbf{the direction of the principle is reversed}.
\item
  \textbf{Quadrature readout --- homodyne detection} {[}10{]}. Reading
  the two quadratures against the phase of a local oscillator is
  standard quantum optics, mathematically identical to our \(x, y\)
  readout. Difference: in optics the LO is an external device supplied
  by the experimenter; here the phase reference is the system's own
  condensate (the proper clock self-selected by entrainment) ---
  \textbf{the universe supplies its own local oscillator}.
\item
  \textbf{The mathematics of the mass readout --- the coherence matrix
  of polarization} {[}11,12{]}. Pauli decomposition of the \(2\times2\)
  coherence matrix = Stokes parameters; \(\det J\ge0\) (equality = full
  polarization) is classical. This paper's contribution is limited to
  the physical identification (depolarization degree = mass\(^2\);
  light-like = fully polarized) and the in-model measurements.
\item
  \textbf{Scale-dependent dimensionality --- the spectral dimension of
  CDT} {[}13{]}. The effective dimension of spacetime runs and decreases
  at short distances. Same direction as our ``ambiguity
  \textasciitilde1/M; non-crystallization in the few-body domain.''
  Difference: theirs is a diffusion dimension; ours is a constructive,
  algebraic dimension gauge (the agreement of rotation axis and cross
  product).
\item
  \textbf{Entrainment --- synchronization theory} {[}14{]}. Entrainment
  into collective rhythms is classical since Huygens. Difference: this
  paper reads entrainment as the \textbf{mechanism by which proper time
  subsists} --- the clock is an order collectively enforced by the
  condensate --- leading to the kinematic classification ``rest is the
  fate of massless perturbations; motion is the privilege of mass.''
\item
  \textbf{Reference frames as physical systems --- quantum reference
  frames} {[}15{]}. Directions and times are defined only relative to
  physical systems (unspeakable information). Difference: there the
  reference frame is given as a resource; here the reference frame
  (clock and axes) is measured to self-generate from the dynamics.
\end{enumerate}

\subsection{15. Outlook --- Open Problems (working guide; not claimed
here)}\label{outlook-open-problems-working-guide-not-claimed-here}

\begin{enumerate}
\def\labelenumi{\arabic{enumi}.}
\tightlist
\item
  \textbf{Matter tracking in intrinsic coordinates only}: track a matter
  packet resisting entrainment using the intrinsic coordinates of
  Section 6, with no external register referenced (completing ``motion =
  the privilege of mass'' intrinsically).
\item
  \textbf{The dynamics of t selection}: place several local condensates
  in a large system and measure how a wide-area time t is selected (the
  collation of clocks = a synchronization network).
\item
  \textbf{The full geometry}: metric via the chord formula
  \(R''=2R\sin(\Delta\theta/2)\), curvature, back-reaction of matter.
\item
  \textbf{The few-body (baryon) domain}: the physics of the domain where
  dimension, time, and existence are ambiguous --- as a separate
  subject.
\item
  \textbf{Attribution of the \(\pm 1\)\% residuals}: resolution or fine
  structure (recorded as low priority).
\item
  \textbf{Stabilization of charge at \(\pm 1\)}: charge (winding)
  distributes over multiple quantized values with no collapse mechanism
  to \(\pm 1\) found (Figure 10). Where the uniqueness of the elementary
  charge comes from is a principal open problem of this series.
  Additionally, the three charge readouts (transmission, relative phase,
  winding) should be measured simultaneously on one state set and
  unified into a single table of monotone relations, sign
  correspondences, conservation, and pair-creation conjugacy.
\item
  \textbf{Continuous sweep of mass degree vs.~entrainment resistance}:
  sweep the non-coherence degree \(\mu=\det\Gamma/T^2\) from 0 to 1 and
  measure lock time, residual detuning, own frequency, and velocity. If
  \(\mu\) and entrainment resistance are related one-to-one, ``mass =
  non-coherence'' advances from an identification to dynamical mass.
  Whether the total \(\sum \det\Gamma\) is conserved or monotone should
  be measured at the same time.
\item
  \textbf{Completeness of the readout algebra}: classify the admissible
  invariant readouts (instantaneous snapshots and finite-window
  correlations; common-phase, permutation, and convention invariant) and
  show that the invariant ring is generated by \(\{x,y,z,\tau,m,q\}\)
  with no independent \(t\) --- turning ``uniquely unreadable'' into a
  no-go theorem.
\item
  \textbf{A statistical arrow of time}: the fundamental dynamics has no
  absolute orientation (Section 11); whether an entropic effective arrow
  emerges from the history of thermalization and autocatalysis is an
  untouched, separate question.
\end{enumerate}

\subsection{16. What This Paper Does Not Claim;
Falsifiability}\label{what-this-paper-does-not-claim-falsifiability}

This paper does not claim that real spacetime is this model.
Falsification conditions: (a) if, from the admissible readouts
(functions of waveform snapshots or finite-window correlations,
invariant under common phase, node permutations, and conventions), a
\textbf{time readout t} independent of \(\tau, m, q\) and independent of
coordinate choices is constructed, the ``uniquely unreadable'' thesis is
rejected --- a form searchable and refutable on the public code. (b) A
clockless perturbation that moves without being entrained. (c) Evidence
that three orthogonal axes exist on the substance side independently of
the readout (e.g., coexistence of an independent complex eigenpair at
equal occupancy in the same condensate). (d) Violation of the 1/M law of
ambiguity. (e) Divergence of mass readouts (i) and (iii) (non-coherence
and entrainment resistance behaving independently).

\textbf{Closing}: the condensate supplied its own position axes and its
own clock --- the axes as the phase decomposition of one complex
coordinate, the clock as the single tick enforced by entrainment. And
exactly one quantity remained unreadable --- coordinate time. Proper
time is physics; coordinate time is a selection --- and this distinction
itself emerged from the readout structure of the closure cone.

\subsection{17. Reproducibility}\label{reproducibility}

We follow the series convention (verification closes within this paper;
originals imported read-only; criteria fixed before execution). The
correspondence table of 8 experiment scripts (including the 4D
translation readout \texttt{run\_translation\_4d\_readout\_v1.py} and
the three-system distribution readout
\texttt{run\_distribution\_readouts\_v1.py}), 2 figure generators, 10
result JSONs plus one npz is canonical in the same folder's experiment
inventory (the experiment inventory file (jikken-ichiran, v1)). All
scripts are deterministic; reproduction agreement with stored JSONs is
machine-checked (performed for v3b/v3c). Experiments whose series were
not stored in early versions (v4) were re-run with identical seeds and
contrast-checked against stored summary values (agreement to six digits)
before being used in figures. The two failed instruments (phase-gradient
gauge; bin quantization) are preserved with their scripts and result
JSONs.

\subsection{References}\label{references}

\textbf{Within the series (provenance labels; the logic of the text
depends on none of them --- all readouts are derived in the text)}

{[}1{]} N. Kihara, \emph{How Inflation Ends Is How Matter Generation
Begins} (2026), Zenodo. Concept DOI: 10.5281/zenodo.21809814 (the
preceding paper: externally supplied position axes and clock;
matter-packet translation; the many-body hair census)

{[}2{]} N. Kihara, \emph{Generation of Fermionic Structure Is Induced,
Autocatalytic, and Pair-Correlated} (2026), Zenodo. Concept DOI:
10.5281/zenodo.21808091 (the 2+1 precursor = the closed-form rotation;
the principle of readout-side classification; the two-body census)

{[}3{]} N. Kihara, \emph{Paper 7: The Temporal Structure of the N-Body
Relational Closed Wave System} (2026), Zenodo. Concept DOI:
10.5281/zenodo.21543070 (discovery of the three-direction metastable
structure)

{[}4{]} N. Kihara, \emph{Paper 8: Causal Separation by Two-Stage Seed
Removal} (2026), Zenodo. Concept DOI: 10.5281/zenodo.21614402 (direct
original of the dynamics code)

{[}5{]} N. Kihara, \emph{Fundamental Axioms of the Anonymous
Equal-Amplitude Composite Wave Model} (2026), Zenodo. Concept DOI:
10.5281/zenodo.21315735 (the zero-square-sum closure --- the axiom of
the cone)

{[}6{]} N. Kihara, \emph{Two-Grammar Decomposition of Interaction}
(Anonymous Two-Channel Closed Wave System Trilogy, Part I) (2026),
Zenodo. Concept DOI: 10.5281/zenodo.21763995 (identification of
transmission with the elementary charge; the charge grammar);
\emph{Future-Phase Position-Acceleration Map in an AB Two-Body Closed
Phase System and the Inverse-Square Law via Harmonic Closure} (2026),
Zenodo. DOI: 10.5281/zenodo.21441081 (provenance of the preceding
paper's external dispersion clock)

{[}7{]} N. Kihara, \emph{The Generative Structure of Fermions} (2026),
Zenodo. Concept DOI: 10.5281/zenodo.21766706 (operational definition of
the closure-completion cost)

\textbf{External}

{[}8{]} D. N. Page, W. K. Wootters, ``Evolution without evolution:
Dynamics described by stationary observables,'' Phys. Rev.~D
\textbf{27}, 2885 (1983)

{[}9{]} C. Rovelli, ``Partial observables,'' Phys. Rev.~D \textbf{65},
124013 (2002)

{[}10{]} U. Leonhardt, \emph{Measuring the Quantum State of Light}
(Cambridge Univ. Press, 1997); H. P. Yuen, V. W. S. Chan, ``Noise in
homodyne and heterodyne detection,'' Opt. Lett. \textbf{8}, 177 (1983)

{[}11{]} G. G. Stokes, ``On the composition and resolution of streams of
polarized light from different sources,'' Trans. Cambridge Philos. Soc.
\textbf{9}, 399 (1852)

{[}12{]} E. Wolf, ``Coherence properties of partially polarized
electromagnetic radiation,'' Nuovo Cimento \textbf{13}, 1165 (1959); M.
Born, E. Wolf, \emph{Principles of Optics}, 7th ed.~(Cambridge Univ.
Press, 1999)

{[}13{]} J. Ambjørn, J. Jurkiewicz, R. Loll, ``Spectral dimension of the
universe,'' Phys. Rev.~Lett. \textbf{95}, 171301 (2005)

{[}14{]} A. Pikovsky, M. Rosenblum, J. Kurths, \emph{Synchronization: A
Universal Concept in Nonlinear Sciences} (Cambridge Univ. Press, 2001)

{[}15{]} S. D. Bartlett, T. Rudolph, R. W. Spekkens, ``Reference frames,
superselection rules, and quantum information,'' Rev.~Mod. Phys.
\textbf{79}, 555 (2007)

\end{document}
